% ============================================================================
%  Informasi Umum Tugas Akhir Mahasiswa -- Teknik Informatika ITERA
%  Templat Baku: Bentuk 1: Penelitian Konvensional
% ============================================================================

% Metadata Judul Naskah (Maksimal 16 kata, format Title Case / UPPERCASE)
\newcommand{\thetitle}{[JUDUL TUGAS AKHIR: OBJEK, METODE, DAN KONTEKS PENELITIAN]}
\newcommand{\thetitleEN}{[THESIS TITLE IN ENGLISH: OBJECT, METHOD, AND RESEARCH CONTEXT]}

% Informasi Identitas Mahasiswa
\newcommand{\theauthor}{[Nama mahasiswa]}
\newcommand{\theauthornim}{[NIM]}
\newcommand{\thetempatlahir}{[Tempat lahir]}
\newcommand{\thetanggallahir}{[Tanggal lahir]}

% Informasi Program Studi & Institusi
\newcommand{\theprodi}{Teknik Informatika}
\newcommand{\thefakultas}{Teknologi Industri}
\newcommand{\theinstitut}{Institut Teknologi Sumatera}
\newcommand{\thelokasi}{Lampung Selatan}
\newcommand{\theyear}{[TAHUN]}
\newcommand{\thekodematkul}{IF25-40033}
\newcommand{\thebentuk}{Bentuk 1: Penelitian Konvensional}

% Informasi Dosen Pembimbing
\newcommand{\namadosbinga}{[Nama dan gelar]}
\newcommand{\nipdosbinga}{NIP. [………………]}
\newcommand{\namadosbingb}{[Nama dan gelar]}
\newcommand{\nipdosbingb}{NIP. [………………]}

% Informasi Dosen Penguji
\newcommand{\namapengujia}{[Nama dan gelar]}
\newcommand{\nippengujia}{NIP. [………………]}
\newcommand{\namapengujib}{[Nama dan gelar]}
\newcommand{\nippengujib}{NIP. [………………]}

% Informasi Koordinator Program Studi (Changelog 1.2.0)
\newcommand{\namakaprodi}{[Nama dan gelar]}
\newcommand{\nipkaprodi}{NIP. [………………]}

% Informasi Tanggal Sidang / Bimbingan
\newcommand{\tanggalsidang}{[tanggal]}

% Makro backward-compatibility untuk template lama
\newcommand{\printnim}{\theauthornim}
\newcommand{\printnamadosbinga}{\namadosbinga}
\newcommand{\printnipdosbinga}{\nipdosbinga}
\newcommand{\printnamadosbingb}{\namadosbingb}
\newcommand{\printnipdosbingb}{\nipdosbingb}
\newcommand{\printnamapengujia}{\namapengujia}
\newcommand{\printnippengujia}{\nippengujia}
\newcommand{\printnamapengujib}{\namapengujib}
\newcommand{\printnippengujib}{\nippengujib}
\newcommand{\printnamakaprodi}{\namakaprodi}
\newcommand{\printnipkaprodi}{\nipkaprodi}
